% ============================================================
% manuscript_sections_simulations_locked_2026-08.tex
% ============================================================
% Locked Simulation 1-4 manuscript text (2026-08), ready to paste into
% Overleaf. manuscript.tex itself only exists in Overleaf (it was removed
% from this git repo in an earlier cleanup commit), so this file is a local
% mirror of the agreed text -- paste each block below into the
% corresponding location in Overleaf. Six insertion points, in order.
% ============================================================


% ------------------------------------------------------------------------
% INSERTION POINT 1
% Replaces everything from "\section{Simulation overview}" down to (but
% not including) "\subsection{Simulation 1...}"
% ------------------------------------------------------------------------

\section{Simulation overview}

The simulations use the model as a controlled theoretical tool. They are not
intended to calibrate a specific clinical population. Instead, they ask what
kinds of patterns can arise when fast symptom dynamics are embedded in a slower
contextual field. Across simulations, the symptom--symptom coupling matrix is
held fixed. What changes is the slow contextual condition, the dynamics of the
slow field, the presence or absence of symptom-to-context feedback, or whether
the slow field is included in the statistical model.

We organize the simulations in four steps. Simulation 1 isolates the effect of
slow contextual baseline: the same fast symptom network is observed under low,
middle, and high values of the slow field. Simulation 2 introduces an acute
perturbation to the slow field and asks how symptoms respond as the slow field
recovers. Simulation 3 adds feedback from sustained symptom activation to the
slow field and asks whether this feedback slows recovery. Simulation 4 turns to
the observation problem and asks what an estimated symptom network recovers when
data are generated by the slow--fast model but the slow field is omitted from
the statistical model.

This order follows the logic of the model. We first show that slow context can
shift symptom burden without changing coupling. We then allow the slow context
to move and recover. We then allow symptoms to feed back into that slow field.
Finally, we ask how the same slow--fast process appears to a researcher who
estimates a symptom-only network.


% ------------------------------------------------------------------------
% INSERTION POINT 2
% Goes AFTER the existing Simulation 1 design/quantities/table.
% Also update the table's burn-in/chains rows to the values below.
% ------------------------------------------------------------------------

\paragraph{Results.}
Changing the slow contextual field shifted symptom burden even though the
symptom--symptom coupling matrix was identical across conditions. In the
low-context condition (\(P=-0.6\)), mean symptom burden was \(M=1.64\), or
\(m=0.18\) of symptoms active. In the middle-context condition (\(P=0\)), mean
burden increased to \(M=2.49\) (\(m=0.28\)). In the high-context condition
(\(P=0.6\)), mean burden increased further to \(M=3.56\) (\(m=0.40\)).

The probability of high symptom burden showed the same pattern. The probability
of having at least half of the symptoms active was \(0.017\) in the low-context
condition, \(0.084\) in the middle-context condition, and \(0.268\) in the
high-context condition. Thus, the same fast symptom network produced different
symptom-burden distributions when embedded in different slow contextual fields.

% Simulation 1 table -- update these three rows:
% \(T_{\mathrm{burn}}\) & Burn-in period & 200 sweeps \\
% \(T_{\mathrm{post}}\) & Post-burn-in samples & 200 sweeps per chain \\
% Chains & Independent trajectories per condition & 200 \\

% FIGURE PLACEHOLDER:
% Figure 2: Simulation 1, burden distributions across P.


% ------------------------------------------------------------------------
% INSERTION POINT 3
% Goes AFTER the existing Simulation 2 design paragraph/table.
% ------------------------------------------------------------------------

\paragraph{Results.}
Before the perturbation, the slow field remained near its baseline
(\(P\approx 0.00\)), and mean symptom burden was about \(M=2.48\) active
symptoms out of 9. At shock onset, the slow field increased to approximately
\(P=1.00\). Symptom burden rose shortly afterward, reaching a peak of about
\(M=4.73\) active symptoms.

By the end of the recovery window, the slow field had returned close to
baseline (\(P\approx 0.06\)), and symptom burden had also largely recovered
(\(M\approx 2.58\)). Thus, an acute perturbation to the slow field produced a
temporary increase in symptom burden without changing the symptom--symptom
coupling matrix. In this simulation, persistence was limited because feedback
from symptoms to the slow field was turned off.

% FIGURE PLACEHOLDER:
% Figure 3: Simulation 2, stress and recovery with feedback off.


% ------------------------------------------------------------------------
% INSERTION POINT 4
% Replaces the ENTIRE existing "\subsection{Simulation 3...}" subsection.
% ------------------------------------------------------------------------

\subsection{Simulation 3: Symptom feedback and persistence}

The third simulation adds feedback from sustained symptom activation to the
slow contextual field. In the previous simulation, the slow field affected
symptoms, but symptoms did not affect the slow field. Here, high symptom
activation can also shift the slow field itself. This captures a simple idea:
symptoms may become more persistent when they start to change the conditions
that maintain them. For example, sustained fatigue, withdrawal, or
concentration problems may worsen work, social, or daily functioning, which in
turn may make later symptom activation more likely.

This simulation is not meant to test for bistability or tipping. Instead, it
asks a simpler question: does symptom-to-context feedback make recovery slower
after a perturbation?

\paragraph{Design.}
Simulation 3 used the same stress-and-recovery design as Simulation 2, but
turned on symptom-to-context feedback. We compared feedback off (\(b=0\)) with
feedback on (\(b=0.50\)). This value was chosen from a small pilot grid because
it produced a visible delay in recovery while all trajectories continued to
decay toward baseline. All other parameters were kept the same as in Simulation
2, including the fast symptom network, the shock magnitude, the mean-reversion
rate, the diffusion scale, and the time step.

Feedback was implemented as
\[
b[\bar m_t-m^\star]_+\Delta t,
\]
where \([x]_+=\max(x,0)\). The smoothed symptom signal \(\bar m_t\) was updated
with an exponential moving average, and \(m^\star\) was set to the mean
active-symptom fraction in the \(P=0\) condition from Simulation 1. Thus, only
sustained symptom activation above the reference level contributed to an upward
shift in the slow field. Below-reference symptom activation was not treated as
an additional protective force.

Because the feedback term uses the positive part, small fluctuations above
\(m^\star\) can create a slight upward pull even before the shock. This effect
was small in the present simulation and did not change the interpretation of
the feedback comparison.

\paragraph{Results.}
Before the shock, both conditions were close to baseline. With feedback off,
the mean slow field during the last 50 pre-shock steps was approximately
\(P=-0.003\), and mean symptom burden was \(M=2.50\). With feedback on, the
corresponding values were \(P=0.013\) and \(M=2.52\).

The immediate post-shock response was similar across conditions, as expected.
In the first 20 post-shock steps, peak symptom burden was \(M=4.34\) with
feedback off and \(M=4.43\) with feedback on. The difference emerged during
recovery. By the end of the recovery window, the slow field had nearly returned
to baseline when feedback was off (\(P=0.05\)), but remained higher when
feedback was on (\(P=0.24\)). Mean symptom burden showed the same pattern,
returning to \(M=2.57\) with feedback off but remaining higher with feedback on
(\(M=2.89\)).

Thus, symptom-to-context feedback made the post-shock state more persistent.
The feedback condition did not show runaway activation or a second stable
high-symptom state. Rather, the same perturbation recovered more slowly when
sustained symptom activation was allowed to feed back into the slow contextual
field.

% FIGURE PLACEHOLDER:
% Figure 4: Simulation 3, feedback off vs on.


% ------------------------------------------------------------------------
% INSERTION POINT 5
% Replaces the ENTIRE existing "\subsection{Simulation 4...}" subsection.
% This is the REPLICATED version (30 reps) -- supersedes any single-run
% Simulation 4 numbers used in earlier drafts.
% ------------------------------------------------------------------------

\subsection{Simulation 4: What estimated networks recover}

The final simulation turns from the data-generating model to the observation
problem. The previous simulations used the slow--fast model to generate symptom
trajectories. Here we ask what happens when data generated by the same model
are analyzed with a symptom-only network model that omits the slow contextual
field.

This simulation is not meant to show that estimated symptom networks are
unusable or that all estimated edges are artifacts. Instead, it asks whether
unmodeled slow context can be absorbed into apparent symptom--symptom coupling.
That is the estimation implication of the slow--fast model.

\paragraph{Design.}
We generated cross-sectional symptom profiles from the revised \(0/1\)
symptom-network model. Each simulated person had one observed symptom profile.
The fast symptom network, thresholds, and context loadings were the same as in
the previous simulations. The true coupling matrix \(\omega\) was therefore
fixed and known.

We compared three arms. In the fixed-context baseline arm, all simulated people
had the same contextual value, \(P=0\). In the pooled symptom-only arm, each
person had a contextual value drawn from
\[
P_i \sim \mathrm{Uniform}(-0.6,0.6),
\]
but the estimated network did not include \(P_i\). In the pooled
context-adjusted arm, the same pooled data were analyzed again, but \(P_i\) was
included as a covariate in the nodewise regressions. Thus, the symptom-only and
context-adjusted arms used the same observations; they differed only in whether
the slow contextual field was modeled.

For each arm, networks were estimated with nodewise logistic regressions and
then symmetrized. We repeated the full procedure across 30 independent
replicates, with 3000 simulated people per arm in each replicate. We compared
estimated global strength, mean absolute recovery error relative to the true
coupling matrix, and the number of phantom edges among symptom pairs with no
true direct coupling.

\paragraph{Results.}
The symptom-only model that omitted the slow field estimated stronger apparent
connectivity than the context-adjusted model. Across 30 replicates, estimated
global strength was highest in the pooled symptom-only arm
(\(M=6.45, SE=0.08\)), lower in the fixed-context baseline arm
(\(M=5.46, SE=0.10\)), and lower still in the context-adjusted arm
(\(M=5.26, SE=0.08\)). The same pattern appeared in recovery error. Mean
absolute error across all edges was \(0.104\) in the symptom-only arm,
\(0.090\) in the fixed-context baseline arm, and \(0.086\) in the
context-adjusted arm.

The difference was especially clear among symptom pairs with no true direct
coupling. Mean absolute error for true-zero edges was \(0.106\) in the
symptom-only arm and \(0.087\) in the context-adjusted arm. The symptom-only
model also produced more phantom edges above 0.10 among truly uncoupled pairs
(\(M=11.83, SE=0.42\)) than the context-adjusted model
(\(M=8.70, SE=0.37\)).

The paired replicate comparisons confirmed that this was a consistent effect.
In all 30 replicates, the symptom-only model estimated higher global strength
than the context-adjusted model. It also had higher recovery error among
true-zero edges in all 30 replicates. On average, omitting the slow field
increased estimated global strength by \(1.19\) (\(SE=0.05\)) and increased
mean absolute error among true-zero edges by \(0.019\) (\(SE=0.001\)).

Thus, when symptom profiles were pooled across different slow-context values,
omitting the slow field made the estimated network look more strongly connected
than the context-adjusted network. This does not mean that estimated edges are
always spurious. It shows that apparent symptom--symptom coupling can partly
reflect unmodeled slow contextual variation.

% FIGURE PLACEHOLDER:
% Figure 5: Simulation 4, replicated naive-adjusted estimation gap.
